\documentclass[10pt,aspectratio=169]{beamer}
\usepackage[portuguese]{babel}
\usepackage{fontspec,unicode-math}
\setmainfont{Noto Sans}
\setsansfont{Noto Sans}
\setmathfont{STIX Math}
\usepackage{amsmath}
\usepackage{xcolor}
\definecolor{InsperRed}{HTML}{a01924}
\definecolor{Ink}{HTML}{232529}
\setbeamercolor{frametitle}{fg=Ink}
\setbeamercolor{structure}{fg=InsperRed}
\setbeamercolor{normal text}{fg=Ink,bg=white}
\setbeamerfont{frametitle}{size=\large,series=\bfseries}
\setbeamerfont{title}{size=\LARGE,series=\bfseries}
\setbeamertemplate{navigation symbols}{}
\setbeamertemplate{itemize item}{\color{InsperRed}$\bullet$}
\setbeamertemplate{footline}{%
  \hbox{\begin{beamercolorbox}[wd=.84\paperwidth,ht=2.7ex,dp=1.2ex,leftskip=.7cm]{normal text}\scriptsize Inteligência Artificial para Decisão · Insper · 2027-1\end{beamercolorbox}%
  \begin{beamercolorbox}[wd=.16\paperwidth,ht=2.7ex,dp=1.2ex,rightskip=.7cm plus1fil]{normal text}\hfill\scriptsize\insertframenumber\, /\, \inserttotalframenumber\end{beamercolorbox}}}
\setbeamersize{text margin left=12mm,text margin right=12mm}
\begin{document}
\begin{frame}[t]{}
\vspace{14mm}\textcolor{InsperRed}{\bfseries AULA 08}\par\vspace{4mm}{\LARGE\bfseries Sistemas conversacionais (chatbots)}\par\vspace{6mm}{\large Finanças}\par\vspace{12mm}{\color{gray}Paulo C. Marques F. · 2027-1}
\end{frame}

\begin{frame}[t]{Objetivos e verificação}
\small
\begin{itemize}
\item Projetar chatbots documentais com histórico, resolução de referências, recuperação contextual e abstenção fundamentada.
\item Resolver manualmente os exemplos do capítulo, identificando hipóteses e consequências.
\item Especificar a aplicação no Codex gráfico e confrontar ao menos um resultado com evidências.
\end{itemize}
\end{frame}

\begin{frame}[t]{Organização do encontro}
\small
\begin{itemize}
\item Apresentação dos fundamentos e resolução de problemas quantitativos no papel.
\item Laboratório supervisionado e testes produzidos por prompts na interface gráfica do Codex, sem terminal.
\item Atividade individual com dados diferentes e depósito no GitHub antes do quiz seguinte.
\end{itemize}
\end{frame}

\begin{frame}[t]{Conversa como estado}
\small
\begin{itemize}
\item Um chatbot documental recebe uma sequência de turnos.
\item A pergunta atual pode depender de referentes introduzidos anteriormente: `Qual é a taxa do Fundo A?'' seguida de `E a liquidez dele?''.
\item A segunda pergunta não é autossuficiente.
\item A consulta reescrita deve incorporar o fundo correto, a classe pertinente e a versão do documento, sem acrescentar informações que o usuário não deu.
\end{itemize}
\end{frame}

\begin{frame}[t]{Conversa como estado — continuação}
\small
\begin{itemize}
\item Podemos representar o estado conversacional por e a transformação de consulta por A consulta é submetida ao retriever.
\item A resposta deve permanecer vinculada aos trechos retornados e não apenas ao estado narrativo da conversa.
\item O histórico ajuda a identificar a pergunta; ele não confirma dados financeiros.
\end{itemize}
\end{frame}

\begin{frame}[t]{Conversa como estado — formulação matemática}
\small
\vspace{3mm}\begin{center}\small\[q'_t=\operatorname{rewrite}(q_t,S_{t-1}).\]\end{center}
\vspace{3mm}\small A consulta é submetida ao retriever.
\end{frame}

\begin{frame}[t]{Referência, ambiguidade e abstenção}
\small
\begin{itemize}
\item Imagine que o diálogo alternou entre Fundos A e B e o usuário pergunta ``Qual deles tem resgate mais rápido?''.
\item Aqui uma comparação é apropriada, mas pode exigir documentos e classes diferentes.
\item Já ``Qual o prazo dele?'' após discutir duas classes sem especificação não possui referente inequívoco.
\item O chatbot deve solicitar esclarecimento.
\end{itemize}
\end{frame}

\begin{frame}[t]{Referência, ambiguidade e abstenção — continuação}
\small
\begin{itemize}
\item O conteúdo deve distinguir taxa de administração, taxa de performance, cotização, liquidação e prazo de resgate.
\item Uma resposta correta para uma classe pode ser incorreta para outra.
\item O campo temporal também é essencial: regulamentos e lâminas são documentos versionados.
\end{itemize}
\end{frame}

\begin{frame}[t]{Fluxo de um chatbot RAG}
\small
\begin{itemize}
\item A aplicação executa recepção da pergunta, atualização de estado, reformulação, recuperação, verificação documental, geração e exibição de fontes.
\item Para cada pergunta, o laboratório deve permitir inspecionar qual consulta foi realmente recuperada.
\item A falha pode nascer na resolução da referência, antes de qualquer busca.
\item Um conjunto de testes deve conter turnos independentes, dependentes do histórico, mudanças explícitas de fundo, mudanças implícitas, perguntas sem documento suficiente e perguntas que exigem abstenção.
\end{itemize}
\end{frame}

\begin{frame}[t]{Fluxo de um chatbot RAG — continuação}
\small
\begin{itemize}
\item A avaliação deve separar correção da consulta reescrita, recuperação de fontes e fundamentação da resposta final.
\end{itemize}
\end{frame}

\begin{frame}[t]{Documentação e responsabilidade}
\small
\begin{itemize}
\item Os documentos didáticos de fundos de investimento serão fictícios, evitando confusão com produtos reais.
\item A aplicação será um assistente de consulta a documentos, e não um consultor de investimentos.
\item Ao usuário deverá ser mostrado o trecho e sua versão, inclusive quando o sistema não consegue responder.
\item O Codex Python SDK poderá integrar a geração, sujeito à validação institucional; a recuperação deve funcionar independentemente do LLM.
\end{itemize}
\end{frame}

\begin{frame}[t]{Exercício para resolver no papel 1}
\small
\small\textbf{8.1 Quem é ``ele''?}\par\vspace{4mm}Histórico: `Qual o prazo de resgate do Fundo Norte?''; E a taxa do Fundo Sul?''; Ela é maior que a do primeiro?''. Identifique a referência de ela'' e de `primeiro''. Reescreva a terceira pergunta como consulta independente sem acrescentar informação.
\end{frame}

\begin{frame}[t]{Exercício para resolver no papel 2}
\small
\small\textbf{8.2 Comparação sem documento}\par\vspace{4mm}O usuário pede `Qual fundo tem menor taxa?'', mas o corpus contém a taxa do Fundo A e não contém documento do Fundo B. É adequado responder `A''? Proponha uma resposta com limitação explícita e a condição que permitiria uma comparação válida.
\end{frame}

\begin{frame}[t]{Laboratório e caso aplicado}
\small
\begin{itemize}
\item Área de aplicação: Finanças.
\item Reformular manualmente três perguntas dependentes do histórico, explicitando fundo, classe e data.
\item Pedir ao Codex interface gráfica de chatbot com histórico, reformulação da consulta, recuperação documental local e citações por resposta.
\item Inspecionar um diálogo em que ``ele'' se refere ao fundo anterior e outro que alterna dois fundos.
\end{itemize}
\end{frame}

\begin{frame}[t]{Prompt de desenvolvimento — exemplo}
\small
\small\textbf{Instrução fornecida ao Codex:}\par\vspace{5mm}Crie um chatbot local para documentos didáticos de três fundos de investimento. Preserve histórico e mostre a consulta reformulada, filtros de fundo/classe/data, trechos recuperados e citações. Respostas sem fonte devem declarar insuficiência; use o Codex SDK somente quando a integração estiver validada.
\end{frame}

\begin{frame}[t]{Critérios de aceitação e validação}
\small
\begin{itemize}
\item Referentes corretos em diálogos com alternância de fundos; números citados na classe/data correspondente; pedido de esclarecimento para identidade incerta; consultas e fontes recuperadas observáveis.
\item Inspecionar argumentos, saídas e rastros das ferramentas quando houver chamadas externas.
\item Comparar uma saída com a resolução manual do exemplo conceitual.
\end{itemize}
\end{frame}

\begin{frame}[t]{Atividade individual em casa}
\small
\begin{itemize}
\item Pacote novo, também fictício, com documentos de dois fundos de investimento e oito perguntas organizadas em três diálogos; incluir alterações entre versões documentais.
\item Resolver as questões de interpretação, com cálculos e evidência documental.
\item Depositar no GitHub os arquivos solicitados; a entrega habilita o quiz, sem valer nota.
\end{itemize}
\end{frame}

\begin{frame}[t]{Preparação para a avaliação escrita}
\small
\begin{itemize}
\item Exercícios do caderno: 8.1 Quem é ``ele''?; 8.2 Comparação sem documento; 8.3 Memória curta, consequência longa.
\item Formalizar o raciocínio e interpretar resultados; não será exigida escrita manual de código.
\item No quiz seguinte serão avaliados conceitos desta aula e da atividade feita em casa.
\end{itemize}
\end{frame}

\end{document}
